% Generated from markdown/01.abstract.md; edit the Markdown source.
Retrieval-augmented generation (RAG) systems rely on retrieval modules to ground large language model (LLM) outputs. LLM-based query expansion enriches retrieval with document-like passages, but evaluations of hybrid retrieval often fuse fixed top-\(L\) prefixes of dense and sparse rankings. Because \(L\) controls cross-channel contributions and ranking access, it can alter measured expansion gains. We therefore evaluate complete-list effectiveness and record per-channel replay stopping depths required to certify the ordered top-\(K\). This changes the design: because both rankings determine the fused result, their query constructions should be coordinated rather than designed independently. We present DESA (Dense Expansion and Sparse Anchoring), which shares generated references across channels but specializes their integration. Orthogonal residual expansion adds new semantic directions to the dense query, whereas score-product anchoring reorders the original sparse support without admitting expansion-only matches. The same references thus play complementary roles: Dense expands; Sparse anchors. Across seven BEIR datasets, DESA improves nDCG@10 and Recall@20 over the unexpanded query by 3.82\% and 2.38\%, while reducing dense and sparse replay stopping depths by 36.90\% and 36.56\%.
